% FRONTEND WEB - diseño de 27 casos (CP-028..CP-054)
% GUÍA EQUIPO: Responsable asignado: Kerly Chicaiza (Frontend Web).
\paragraph{Módulo Web --- Tienda en línea y panel de rutas (27 casos: CP-028 a CP-054).}

\begin{disenio}
\caso{CP-028}{\textbf{Historial unificado Mis pedidos con compras SHOP y ADMIN propias.} NR-01, NR-07, NP-19 (C05.1, Alta)}{Sesión web de CLIENTE-A; una orden de origen SHOP y una de origen ADMIN del mismo titular.}{1) Abrir Mis pedidos en web. 2) Verificar listado vía GET /\allowbreak{}api/\allowbreak{}v1/\allowbreak{}shop/\allowbreak{}orders. 3) Abrir el detalle de ambas órdenes.}{La lista unificada muestra solo las dos compras propias de CLIENTE-A con su origen (SHOP/\allowbreak{}ADMIN) y estado de seguimiento.}{Coinciden identidades, orígenes y seguimiento entre UI y API; no aparecen órdenes de terceros ni invitadas.}
\caso{CP-029}{\textbf{Bloqueo en lista y detalle ante pedido ajeno o invitado coincidente.} NR-01, NR-07, NP-19 (C05.2, Alta)}{Sesión web de CLIENTE-A; ID de pedido de CLIENTE-B y pedido invitado con teléfono coincidente.}{1) Revisar Mis pedidos de CLIENTE-A. 2) Intentar abrir en web y GET /\allowbreak{}api/\allowbreak{}v1/\allowbreak{}shop/\allowbreak{}orders/\allowbreak{}:\allowbreak{}id el pedido ajeno y el invitado.}{El pedido invitado no se vincula por teléfono en la lista; ambos accesos por ID responden 404 ORDER\_NOT\_FOUND sin datos.}{Sin fuga de información en UI ni API; el código 404 ORDER\_NOT\_FOUND bloquea detalle ajeno e invitado.}
\caso{CP-030}{\textbf{Pedido sin PurchaseReceipt y sin importes inventados.} NR-05, NR-08, NP-03 (C08.1, Media)}{Sesión web de CLIENTE-A; ID de pedido previo propio con receiptAvailable=\allowbreak{}false.}{1) Consultar lista y detalle del pedido en web y GET /\allowbreak{}api/\allowbreak{}v1/\allowbreak{}shop/\allowbreak{}orders. 2) Solicitar /\allowbreak{}api/\allowbreak{}v1/\allowbreak{}shop/\allowbreak{}orders/\allowbreak{}:\allowbreak{}id/\allowbreak{}receipt.}{Lista y detalle marcan receiptAvailable=\allowbreak{}false e informan comprobante no disponible; /\allowbreak{}receipt responde 404 RECEIPT\_UNAVAILABLE.}{No se reconstruye comprobante ni se inventan importes o impuestos en UI ni en API; datos almacenados intactos.}
\caso{CP-031}{\textbf{Cobertura tienda: punto interior al polígono aceptado en web.} NR-11, NR-12, NR-13, NP-05 (C09.1, Alta)}{Sesión CLIENTE-A; coordenadas sintéticas dentro del polígono de cobertura vigente; carrito Shop válido.}{1) Seleccionar punto interior en el mapa de checkout web. 2) Verificar POST /\allowbreak{}api/\allowbreak{}v1/\allowbreak{}store-coverage/\allowbreak{}evaluate. 3) Completar compra.}{La evaluación retorna position=\allowbreak{}INTERIOR y accepted=\allowbreak{}true; la interfaz permite avanzar y confirmar el pedido Shop.}{El punto interior queda aceptado sin bloqueos y la orden se registra con sus coordenadas válidas.}
\caso{CP-032}{\textbf{Cobertura tienda: punto exactamente en el borde del polígono.} NR-11, NR-12, NR-13, NP-05 (C09.2, Alta)}{Sesión CLIENTE-A; coordenadas exactamente sobre el borde del polígono de StoreCoverage.}{1) Ingresar punto en el borde en checkout web. 2) Evaluar vía POST /\allowbreak{}api/\allowbreak{}v1/\allowbreak{}store-coverage/\allowbreak{}evaluate. 3) Confirmar si permite compra.}{La evaluación devuelve position=\allowbreak{}BOUNDARY y accepted=\allowbreak{}true, habilitando la compra en la web.}{El punto en el borde es aceptado consistentemente en UI y API según la regla geométrica vigente.}
\caso{CP-033}{\textbf{Cobertura tienda: punto exterior rechazado con 422 en web.} NR-11, NR-12, NR-13, NP-05 (C09.3, Alta)}{Sesión CLIENTE-A; coordenadas fuera del polígono de cobertura; producto disponible en carrito.}{1) Seleccionar punto exterior en checkout web. 2) Consultar evaluate e intentar confirmar compra. 3) Verificar stock y órdenes.}{Evaluate retorna EXTERIOR y accepted=\allowbreak{}false; el intento de compra se rechaza con 422 DELIVERY\_OUTSIDE\_COVERAGE y aviso en UI.}{No se crea pedido ni comprobante, y el stock y saldo simulado del cliente permanecen sin cambios.}
\caso{CP-034}{\textbf{Cambio de cobertura A a B conserva pedidos ya aceptados.} NR-11, NR-12, NR-13, NP-05 (C09.4, Alta)}{Polígono A inicial con pedido aceptado en punto P; polígono B nuevo que excluye al punto P.}{1) Como ADMIN actualizar polígono a B en panel web (PUT /\allowbreak{}api/\allowbreak{}v1/\allowbreak{}admin/\allowbreak{}store-coverage). 2) Consultar el pedido previo.}{El cambio de cobertura se guarda, pero el pedido aceptado en P conserva su estado, dirección y datos sin cancelación.}{Ningún pedido previo aceptado es cancelado, alterado ni eliminado automáticamente tras el cambio de polígono.}
\caso{CP-035}{\textbf{Compra nueva tras cambio de cobertura aplica área B vigente.} NR-11, NR-12, NR-13, NP-05 (C09.5, Alta)}{Sesión CLIENTE-A; punto P válido en área A pero exterior al área B vigente; carrito válido.}{1) Intentar nueva compra web para punto P bajo área B. 2) Verificar respuesta y estado de carrito/\allowbreak{}pedidos.}{La compra nueva se rechaza en UI y API con 422 DELIVERY\_OUTSIDE\_COVERAGE al aplicar el polígono B vigente.}{Sin pedido, comprobante ni cobro nuevos; el pedido ya aceptado conserva su dirección original.}
\caso{CP-036}{\textbf{Horario Shop límite inicial 07:\allowbreak{}00 America/\allowbreak{}Guayaquil aceptado.} NR-14, NP-06 (C10.1, Alta)}{Sesión CLIENTE-A; fecha futura válida con hora 07:\allowbreak{}00 (-05:\allowbreak{}00 America/\allowbreak{}Guayaquil); pago DEMO\_APPROVED.}{1) Seleccionar 07:\allowbreak{}00 en checkout web. 2) Confirmar compra vía POST /\allowbreak{}api/\allowbreak{}v1/\allowbreak{}shop/\allowbreak{}checkout. 3) Consultar pedido.}{El horario 07:\allowbreak{}00 se acepta en UI y backend; la orden se crea con la hora programada correspondiente.}{HTTP 201 en checkout, pedido y comprobante generados y hora almacenada coherente con 07:\allowbreak{}00 America/\allowbreak{}Guayaquil.}
\caso{CP-037}{\textbf{Horario Shop límite final 17:\allowbreak{}00 America/\allowbreak{}Guayaquil aceptado.} NR-14, NP-06 (C10.2, Alta)}{Sesión CLIENTE-A; fecha futura válida con hora 17:\allowbreak{}00 (-05:\allowbreak{}00 America/\allowbreak{}Guayaquil); carrito y cobertura válidos.}{1) Seleccionar 17:\allowbreak{}00 en checkout web. 2) Confirmar la compra y consultar el detalle del pedido.}{El horario límite 17:\allowbreak{}00 es aceptado en web y API, creando la orden con entrega programada a las 17:\allowbreak{}00.}{Compra confirmada (201) sin rechazos indebidos en el extremo superior 17:\allowbreak{}00.}
\caso{CP-038}{\textbf{Horario Shop inválido 06:\allowbreak{}59 rechazado sin orden ni cobro.} NR-14, NP-06 (C10.3, Alta)}{Sesión CLIENTE-A; fecha futura válida con hora 06:\allowbreak{}59 (-05:\allowbreak{}00 America/\allowbreak{}Guayaquil); tarjeta DEMO\_APPROVED.}{1) Intentar confirmar compra para las 06:\allowbreak{}59 en web o solicitud checkout. 2) Verificar stock, saldo y órdenes.}{Se rechaza con 422 DELIVERY\_OUTSIDE\_HOURS y mensaje claro en UI; no se genera orden ni comprobante.}{Stock y saldo simulado intactos tras el rechazo; cero órdenes creadas para 06:\allowbreak{}59.}
\caso{CP-039}{\textbf{Horario Shop inválido 17:\allowbreak{}01 rechazado sin orden ni cobro.} NR-14, NP-06 (C10.4, Alta)}{Sesión CLIENTE-A; fecha futura válida con hora 17:\allowbreak{}01 (-05:\allowbreak{}00 America/\allowbreak{}Guayaquil); carrito válido.}{1) Intentar confirmar compra para las 17:\allowbreak{}01 en checkout web. 2) Verificar órdenes, stock y saldo.}{El servidor responde 422 DELIVERY\_OUTSIDE\_HOURS y la web muestra error sin crear pedido ni realizar cobro.}{Sin orden ni comprobante emitidos; inventario y saldo del cliente permanecen sin variación.}
\caso{CP-040}{\textbf{Cambio de horario en shop-settings genera 409 al confirmar.} NR-14, NP-06 (C10.9, Alta)}{Cotización Shop abierta con horario 07:\allowbreak{}00-17:\allowbreak{}00; ADMIN cambia ventana vía PATCH /\allowbreak{}api/\allowbreak{}v1/\allowbreak{}admin/\allowbreak{}shop-settings.}{1) Cotizar en web como CLIENTE-A. 2) Como ADMIN actualizar horario. 3) Confirmar compra pendiente en web.}{El checkout responde 409 SHOP\_SETTINGS\_CHANGED; la web informa el cambio y exige recotizar sin confirmar con reglas previas.}{No se crea orden ni cobro bajo el horario obsoleto; tras recotizar aplican las condiciones vigentes.}
\caso{CP-041}{\textbf{Verificación UI web de aviso de hora de entrega aproximada.} NR-14, NP-06 (C10.10, Alta)}{Sesión CLIENTE-A; pantalla de selección de horario en checkout Shop y detalle de pedido creado.}{1) Inspeccionar textos de horario en checkout web. 2) Completar compra y revisar textos en detalle del pedido.}{La interfaz informa claramente que la hora de entrega es aproximada y no promete exactitud ni margen de +-30 min.}{El aviso visible es explícito sobre carácter aproximado y no garantiza ventana +-30 min (validación puramente UI).}
\caso{CP-042}{\textbf{Propuesta de rutas en panel ADMIN es solo lectura sin persistir.} NR-20, NP-07, NP-13 (C14.1, Alta)}{Sesión ADMIN en panel web; pedidos elegibles de hoy (isToday=\allowbreak{}true) y otro día (isToday=\allowbreak{}false).}{1) Solicitar propuesta en panel (GET /\allowbreak{}api/\allowbreak{}v1/\allowbreak{}route-proposals) para hoy y fecha futura. 2) Consultar rutas persistidas.}{Devuelve propuestas temporales con isToday/\allowbreak{}canConfirm correctos; consultar rutas persistidas muestra cero rutas creadas.}{Ninguna ruta ni reserva se persiste en BD al generar propuestas; canConfirm=\allowbreak{}false en fechas distintas de hoy.}
\caso{CP-043}{\textbf{Confirmación de borrador DRAFT en hoy persiste ruta sin conductor.} NR-20, NP-07, NP-13 (C14.2, Alta)}{Sesión ADMIN; propuesta del día actual con canConfirm=\allowbreak{}true y pedidos elegibles.}{1) Confirmar propuesta en panel ADMIN (POST /\allowbreak{}api/\allowbreak{}v1/\allowbreak{}routes). 2) Consultar listado de rutas y reservas de pedidos.}{Se persiste una ruta en estado DRAFT sin conductor asignado (driverId=\allowbreak{}null) y visible en el panel.}{Estado DRAFT confirmado sin duplicar rutas ni asignar conductor automáticamente; pedidos reclamados por la ruta.}
\caso{CP-044}{\textbf{Asignación separada de ruta DRAFT a conductor disponible.} NR-20, NP-07, NP-13 (C14.3, Alta)}{Sesión ADMIN; ID de ruta DRAFT y REPARTIDOR-A sin ausencias ni reservas activas.}{1) Desde panel ADMIN asignar REPARTIDOR-A (POST /\allowbreak{}api/\allowbreak{}v1/\allowbreak{}routes/\allowbreak{}:\allowbreak{}id/\allowbreak{}assign). 2) Consultar ruta y disponibilidad.}{La ruta pasa a estado ASSIGNED con REPARTIDOR-A y crea reserva activa del conductor en el sistema.}{Transición DRAFT\ensuremath{\rightarrow}ASSIGNED persistida como acción separada; reserva creada y reflejada en disponibilidad.}
\caso{CP-045}{\textbf{Rechazo de asignación a conductor no disponible o ausente (409).} NR-20, NP-07, NP-13 (C14.4, Alta)}{Sesión ADMIN; ruta DRAFT sin conductor y REPARTIDOR-B con ausencia vigente o reserva previa.}{1) Intentar asignar REPARTIDOR-B a la ruta DRAFT en panel ADMIN. 2) Consultar rutas sin conductor y reservas.}{La asignación se rechaza con HTTP 409 (DRIVER\_ABSENT o conflicto de reserva) y aviso claro en UI.}{Sin asignación parcial ni reserva indebida; la ruta permanece en DRAFT sin conductor y visible en el panel.}
\caso{CP-046}{\textbf{Motor vial OSRM indisponible muestra aviso sin rutas en línea recta.} NR-19, NR-21, NP-09 (C16.1, Alta)}{Sesión ADMIN; pedidos elegibles de hoy con motor vial OSRM detenido o inaccesible.}{1) Solicitar propuestas en panel ADMIN (GET /\allowbreak{}api/\allowbreak{}v1/\allowbreak{}route-proposals). 2) Observar aviso UI y respuesta HTTP.}{El servidor responde 503 ROAD\_ENGINE\_UNAVAILABLE (o 502) y el panel muestra aviso claro sin generar rutas en línea recta.}{Cero propuestas ficticias o trazados en línea recta presentados como resultado vial; nada se persiste.}
\caso{CP-047}{\textbf{Preservación de pedidos pendientes y rutas tras falla OSRM.} NR-19, NR-21, NP-09 (C16.2, Alta)}{Sesión ADMIN; estado previo de pedidos PENDING y ruta existente antes y después del error 503/\allowbreak{}502.}{1) Consultar listado de pedidos pendientes y rutas existentes en panel tras el rechazo 503/\allowbreak{}502.}{Los pedidos siguen en PENDING y las rutas previamente confirmadas o asignadas se conservan intactas.}{Comparación antes/\allowbreak{}después confirma ausencia de mutaciones, cancelaciones o reservas parciales tras el error vial.}
\caso{CP-048}{\textbf{Recuperación tras indisponibilidad OSRM genera propuestas válidas.} NR-19, NR-21, NP-09 (C16.3, Alta)}{Sesión ADMIN; servicio OSRM nuevamente operativo y pedidos pendientes del día.}{1) Verificar estado en panel/\allowbreak{}road-engine. 2) Solicitar nuevamente propuestas desde panel ADMIN.}{El panel genera propuestas viales válidas con rutas, distancias y duraciones estimadas normalmente.}{HTTP 200 en propuestas tras recuperación; resultados coherentes con el motor vial sin errores residuales.}
\caso{CP-049}{\textbf{Origen y retorno común en tienda única reflejado en propuesta y mapa.} NR-11, NR-19, NP-08 (C17.1, Media)}{Sesión ADMIN; pedidos del día con destino válido dentro de cobertura y tienda única configurada.}{1) Solicitar propuesta con geometría en panel ADMIN. 2) Verificar punto de inicio y fin en JSON y mapa.}{La propuesta y el mapa web parten y regresan a la ubicación única de StoreCoverage.\allowbreak{}location, sin depósitos por zona.}{Coordenadas de origen y retorno coinciden exactamente con StoreCoverage.\allowbreak{}location configurada.}
\caso{CP-050}{\textbf{Paradas y trazado vial en panel ADMIN coinciden con motor vial.} NR-11, NR-19, NP-08 (C17.2, Media)}{Sesión ADMIN; propuesta con 2 a 4 paradas en horarios compatibles y trazado geométrico OSRM.}{1) Abrir mapa de la propuesta en panel ADMIN. 2) Comparar secuencia de paradas y polilínea con respuesta API.}{El mapa muestra las paradas en el orden cronológico devuelto y sigue el trazado vial OSRM, no rectas entre puntos.}{Orden de paradas y geometría renderizada corresponden fielmente con los datos viales recibidos del backend.}
\caso{CP-051}{\textbf{Distancias y duraciones mostradas como estimadas sin tráfico en vivo.} NR-11, NR-19, NP-08 (C17.3, Media)}{Sesión ADMIN; metadatos roadMetrics (estimate=\allowbreak{}true, trafficAware=\allowbreak{}false) y resumen mostrado en UI.}{1) Revisar métricas e indicadores en panel de propuestas. 2) Contrastar textos visibles con roadMetrics del API.}{La UI presenta distancias y duraciones como estimadas (estimate=\allowbreak{}true, trafficAware=\allowbreak{}false) sin afirmar ETA ni tráfico vivo.}{Ninguna etiqueta web promete hora exacta de llegada, tráfico en tiempo real ni solución óptima global.}
\caso{CP-052}{\textbf{Regresión Shop: compra válida con DEMO\_APPROVED crea pedido y recibo.} NR-01, NR-05, NR-10 (C27.1, Alta)}{Sesión CLIENTE-A; producto disponible, cobertura y horario válidos, tarjeta simulada DEMO\_APPROVED.}{1) Añadir producto desde catálogo Shop web. 2) Completar checkout con DEMO\_APPROVED. 3) Abrir pedido y comprobante.}{Se crea el pedido y su comprobante simulado; el stock y el saldo demo se descuentan exactamente una vez.}{Comprobante coincide con precios e IVA del servidor; sin exponer número completo ni CVV de tarjeta en evidencias.}
\caso{CP-053}{\textbf{Regresión Shop: carrito mixto con agotado rechazado sin efectos.} NR-01, NR-05, NR-10 (C27.2, Alta)}{Sesión CLIENTE-A; carrito con 1 producto disponible y 1 agotado; stock y saldo inicial medidos.}{1) Intentar cotizar o confirmar carrito mixto en Shop web/\allowbreak{}API. 2) Verificar órdenes, stock disponible y saldo.}{El catálogo marca el agotado como no disponible y el checkout rechaza el carrito mixto sin crear orden parcial.}{Stock del producto disponible y saldo demo quedan idénticos al estado previo; ningún pedido o recibo creado.}
\caso{CP-054}{\textbf{Regresión Shop: precio o total manipulado en cliente no es aceptado.} NR-01, NR-05, NR-10 (C27.3, Alta)}{Sesión CLIENTE-A; payload hacia POST /\allowbreak{}api/\allowbreak{}v1/\allowbreak{}shop/\allowbreak{}checkout con campos price/\allowbreak{}total alterados a valor menor.}{1) Enviar solicitud manipulada interceptada desde cliente. 2) Inspeccionar respuesta HTTP, orden y cobro.}{El servidor rechaza campos no permitidos (400/\allowbreak{}409) o aplica únicamente sus precios reales; nunca cobra el monto alterado.}{Si rechaza no deja orden ni débito; si acepta solo cobra e imprime el importe calculado por el servidor.}
\end{disenio}
